\section{Introduction}

\subsection{Background of the Study}

Due to digitalization and the change of behavior of consumer in QSR industry for the last few years \citep{govender2024digital} have been transformed a lot. Perhaps the biggest change has been from third-party delivery platforms, which have transformed customer relationships with restaurants \citep{meemken2022research}. Mobile applications such as Wolt and Foodora have become core players in the food service ecosystem and, through them, customers can place orders for meals over this application quickly and receive on-demand delivery within minutes \citep{simonovits2022uberisation, tuomi2023riding}.

Urbanization, rising smartphone usage and a preference for convenience have contributed to the proliferation of these platforms. This demand has increased consumer expectations of speed, reliability and digitally accessible food delivery services -- particularly in urban markets \citep{siddiqi2024investigating}. Consequently, QSR businesses found themselves heavily reliant on third-party platforms to serve as a primary or secondary sales channel. The new landscape has not only broadened market access but also added complexity to restaurants both operationally and strategically \citep{sun2022too}.

One of the key ramifications of this platform-based model is the massive data it creates. Transactions processed on a platform, such as Wolt and Foodora, generate extensive data logs with information about the time of orders, order values, product selection, time at which goods were delivered and service provider fees per platform \citep{nayak2024future}. In this process, you end up with historical datasets that form insights on customer behavior patterns, seasonal demand variations and comparative performance of the platform. That data, theoretically, is a valuable basis for predictive analytics and decision making.

Nonetheless, this data is readily available but many QSRs do not leverage it to its full potential \citep{almeida2023time,nayak2022restaurant} Data is in fact, usually used for the purpose of describing events that already happened (like looking back at previous sales or coming up with simple reports). Although these descriptive insights are informative in that they provide information about what happened, they do not predict future outcomes or how to intervene to get better outcomes \citep{ma2020making, roy2022complete}. Between making data available and providing intelligence to act upon there exists an insurmountable gulf.




This rising standard is reflected through the same increasing competitive nature of the QSR industry \citep{altenburger2019artificial}. Between these managers are various delivery platforms that restaurants need to work with at once, each having different commissions, customer bases and promotional dynamics \citep{feldman2023managing}. For example, deciding how to spread budgets by channel, timing of sales/discounting or responding to demand volatility has a direct effect on revenue and bottom line. It is sub-optimal to simply be decisioning these based on gut feels or historical averages \citep{cai2023marketing}.

Data-driven decision making has emerged as an important capability in this context. Businesses are progressively also expected to shift from descriptive analytics towards predictive and prescriptive \citep{berkhout2024revisiting}. On the one hand, predictive analytics allows organizations to make predictions about future demand based on historical data; while prescriptive analytics takes it further and recommends which actions should be taken to achieve objectives (e.g., maximize profit or improve service efficiency) rather than reporting only future benefits [4].

For QSRs existing on third-party platforms, the combination with these analytic methods has great potential \citep{li2024criticality}. Through the systematic analysis of platform data, businesses can spot demand patterns with respect to time and platforms, predict future sales trends well in advance, as well as take note to undertake preemptive actions \citep{wu2025does}. For example, PR frameworks can help in generating accurate demand forecasts leading to better staffing, inventory planning and promotional strategies whilst optimization models can be used for platform level decision making improvements \citep{shen2025jd}.


However, there is still a shortage of frameworks that combine the challenges of data integration, forecasting and decision optimization in one coherent framework \citep{khan2025model}. The tools or expertise required to derive actionable insights from raw data that is useful for managerial decisions are hard to come by for many businesses \citep{zhang2025timeraf}. Hence, this study is driven by the purpose of establishing a data-driven framework that combines both availability and usability of data in decision-making within the context of QSRs, which is an essential practice for gainful business implementation as well as an addition to the existing field of business analytics.

\subsection{Problem Statement}

With a more competitive third-party delivery marketplace, and the transaction data of sales widely available with merchants \citep{ogunmokun2025role}, coming up with a tactical plan from the ground-level (sales data) can be much harder for many businesses \citep{taylor2017new}. For HungryHub Oy, big amounts of transactional data is continuously produced via platforms like Wolt and Foodora \citep{mbare2023psychosocial}. So we gather data on order trends, customer preferences, how the platform is performing.


But the current utility of this data is still frozen in time. In the practice it is mainly identifiable as descriptive analysis like using historical trends of sales for reporting purposes to create simple reports \citep{wolniak2023concept}. Though this gives a glimpse into how things were, it does not help plan for the future. Consequently, important operational and strategic choices—e.g., allocation of sales on platform-level, promotion planning and demand forecasting—are not sufficiently backed by systematic analytic approaches \citep{groene2024introduction}.


These limited use of advanced data creates a massive divide between the data available and decision making effectiveness. So, HungryHub is unable to provide a good demand forecasting for future across multiple platforms without predictive models. Moreover, the company is missing prescriptive tools which provide possible overall optimal decisions to maximize profit from forecast. Rationality might be replaced with intuition or naive assumptions in decision-making, which can result in substantive inefficiencies and missed opportunities culminating in unsatisfactory financial outcomes.


This constraint lessens the companys ability to preemptively react to fluctuations in demand and optimize their operational functionality across delivery channels in a highly competitive QSR space \citep{racine2022college}. Hence, the main issue tackled in this work is the lack of a unified analytical framework incorporating data preparation, demand forecasting, and decision optimization.

\subsection{Research Objectives}

The main aim of this research is creation of a prescriptive decision support system to optimize the sales of third party delivery platform. In particular, the study will seek to combine sales data on websites like Wolt and Foodora into a single analytical dataset that will allow performing a thorough analysis. 

Based on this, the research aims at constructing and testing the models of forecasting that will have the capacity to predict the demand at the platform level. Moreover, the study aims at converting these demand predictions into profit-driven decision makers by developing an optimization model. Lastly, the study will evaluate the possible business consequences of the proposed framework through simulation or back-testing, to determine the effectiveness of the proposed framework in enhancing decision-making and the overall sales performance.

\subsection{Research Questions}

The study addresses the following research questions:

\begin{itemize}
    \item \textbf{RQ1:} How can sales data from third-party delivery platforms be structured into a unified analytical dataset for decision-making?
    
    \item \textbf{RQ2:} Which forecasting approach provides the most accurate predictions of platform-level demand?
    
    \item \textbf{RQ3:} How can demand forecasts be transformed into profit-oriented decision recommendations?
    
    \item \textbf{RQ4:} What is the potential business impact of a prescriptive decision-support framework in optimizing platform sales performance?
\end{itemize}

\subsection{Scope and Limitations}

This research project aims at developing a data-driven decision support model of third party delivery sales at HungryHub Oy. Only sales data of two platforms, Wolt and Foodora, is analyzed. The research primarily investigates platform-level performance such as sales trend and demand prediction.

The study involves the cleaning and integration of the data, the formulation of forecasting models and the formulation of a simple profit-based decision model. The outcomes are measured on the basis of simulation or back-testing of the possible impact on business.

But there exist certain limitations. Other channels of sales (dine-in or takeaway) are not part of the study. It also does not discuss such operational aspects as kitchen management, staffing or delivery logistics. Moreover, the framework is not used as an active system, but merely tested with historical data. These constraints allow maintaining the research narrow and within the scope of a Master’s thesis.

\subsection{Structure of the Thesis}
The structure of this thesis is summarized in Table~\ref{tab:thesis_structure}.

\begin{table}[ht]
\centering
\renewcommand{\arraystretch}{1.5}
\begin{tabular}{|c|p{10cm}|}
\hline
\textbf{Chapter} & \textbf{Description} \\ \hline

\textbf{Chapter 1} &
Introduction to the study, including background, problem statement,
research objectives, research questions, and scope. \\ \hline

\textbf{Chapter 2} &
Theoretical background covering business analytics, demand forecasting,
decision-support systems, and Design Science Research. \\ \hline

\textbf{Chapter 3} &
Research methodology, including research approach, data collection,
and analysis methods. \\ \hline

\textbf{Chapter 4} &
Data preparation and development of forecasting models. \\ \hline

\textbf{Chapter 5} &
Design of the decision-support framework and evaluation using simulation
or back-testing. \\ \hline

\textbf{Chapter 6} &
Discussion of results, key findings, and practical recommendations. \\ \hline

\textbf{Chapter 7} &
Conclusion, including summary, limitations, and suggestions for future research. \\ \hline

\end{tabular}
\caption{Structure of the Thesis}
\label{tab:thesis_structure}
\end{table}